\section{Datasets Statistics}
\label{app:a}
We conduct evaluations on six widely used RAG benchmarks from the FlashRAG toolkit~\cite{FlashRAG}, covering both single-hop and multi-hop question answering tasks:
\begin{itemize}
    \item \textbf{Natural Questions (NQ)~\cite{nq}.} Real user questions from Google Search paired with Wikipedia passages/answers; commonly used for open-domain, mostly single-hop QA.
    \item \textbf{PopQA~\cite{popqa}.} Popular-knowledge question set designed for retrieval-based QA, emphasizing factual queries where the answer must be grounded in retrieved evidence.
    \item \textbf{TriviaQA~\cite{triviaqa}.} Trivia-style questions with evidence documents (often web/Wikipedia); used for open-domain factual QA and long-context evidence matching.
    \item \textbf{HotpotQA~\cite{hotpotqa}.} Multi-hop QA requiring reasoning over multiple supporting Wikipedia passages; includes labeled supporting facts.
    \item \textbf{Musique~\cite{musique}} Multi-hop QA benchmark built to test compositional reasoning across several pieces of evidence, often with more challenging, structured multi-step requirements.
    \item \textbf{2WikiMultiHopQA (2Wiki)~\cite{2wiki}} Multi-hop QA constructed from Wikipedia that typically requires linking two (or more) pages to reach the answer, focusing on cross-article reasoning.
\end{itemize}
The detailed dataset statistics are in Table~\ref{tab:dataset}. For RL-based agentic systems, we randomly sample 5000 data from the Train set of HotpotQA and NQ. For the test set of RAGSearch, we adopt the full original Dev or Test set. In detail, for NQ, PopQA and TriviaQA, we choose the Test set; For HotpotQA, Musique, 2Wiki, we choose the Dev set.
\begin{table*}[!h]
\caption{Dataset Statistics}
\label{tab:dataset}
\begin{tabular}{c|c|c|ccc}
\hline
Dataset  & Task         & Knowledge Source & \#Train                                               & \#Dev                                                & \#Test                                           \\ \hline
NQ       & General QA   & Wiki             & 79,168                                                & 8,757                                                & 3,610                                            \\
PopQA    & General QA   & Wiki             & -                                                     & -                                                    & 14,267                                           \\
TriviaQA & General QA   & Wiki \& Web   & 78,785   & 8,837 & 11,313                                           \\
HotpotQA & Multi-hop QA & Wiki             & 90,447                                                & 7,405                                                & -                                                \\
Musique  & Multi-hop QA & Wiki             & 19,938                                                &  2,417 &  - \\
2Wiki    & Multi-hop QA & Wiki             & 15,000                                                & 12,576                                               & -                                                \\ \hline
\end{tabular}
\end{table*}
\section{Implementation Details}
\label{app:b}
To conduct a comprehensive evaluation, we include two different categories of agentic systems for assessment. Here is brief introduction to all the method involved.
\paragraph{\textbf{Training-free Workflows:}}
These approaches do not train an explicit control policy; rather, they use structured prompts and heuristic rules to steer multi-step retrieval and reasoning at inference time. Specifically, we evaluate;
\begin{itemize}
    \item Search-o1~\cite{li2025searcho1}: Augments large reasoning models with an agentic RAG workflow and a Reason-in-Documents module that refines retrieved evidence before integration, enabling dynamic, noise-reduced knowledge retrieval to improve reliability on complex reasoning tasks and open-domain QA. Our implementation is based on \url{https://github.com/RUC-NLPIR/Search-o1}
    \item GraphSearch~\cite{yang2025graphsearch}: An agentic deep-search workflow for GraphRAG that performs multi-turn, modular retrieval with dual-channel querying over both text chunks (semantic) and structural graphs (relational), consistently improving multi-hop RAG accuracy and generation quality over traditional GraphRAG retrieval. Our implementation is based on \url{https://github.com/DataArcTech/GraphSearch}
\end{itemize}
\paragraph{\textbf{Trained Search Agent:}}
These approaches use an RL policy to optimize the agent’s search and reasoning behavior, typically adopting GRPO as the reinforcement learning algorithm. Specifically, we choose:
\begin{itemize}
    \item Search-R1~\cite{jin2025searchr1}: A RL-based retrieval-augmented reasoning framework that trains LLMs to autonomously generate multi-turn search queries during step-by-step reasoning, using retrieved-token masking and an outcome-based reward to improve QA performance over standard RAG baselines. Our implementation is based on \url{https://github.com/PeterGriffinJin/Search-R1}
    \item Graph-R1~\cite{luo2025graphr1}: An agentic GraphRAG framework trained end-to-end with reinforcement learning that builds lightweight knowledge hypergraphs and performs multi-turn retrieval as an agent–environment interaction. Our implementation is based on \url{https://github.com/LHRLAB/Graph-R1}
\end{itemize}
For all the RL-based methods, we pre-train the model with GRPO for 3 epochs. We set the batch size as 32, max search turn as 5.

For retrieve backends, we mainly adopt 5 representative GraphRAGs:
\begin{itemize}
    \item \textbf{HypergraphRAG}~\cite{luo2025hypergraphrag}: A hypergraph-based RAG framework that represents real-world n-ary facts using hyperedges and integrates hypergraph construction, retrieval, and generation. Our implementation is based on \url{https://github.com/LHRLAB/Graph-R1}
    \item \textbf{HippoRAG2}~\cite{gutiérrez2025hipporag2}:A memory-inspired RAG framework that extends HippoRAG’s Personalized PageRank retrieval with deeper passage integration and stronger online LLM usage, improving factual, sense-making, and associative memory. Our implementation is based on \url{https://github.com/OSU-NLP-Group/HippoRAG}
    \item \textbf{LinearRAG}~\cite{zhuang2025linearrag}: An efficient GraphRAG framework that avoids noisy, costly relation extraction by building a lightweight relation-free hierarchical “Tri-Graph” (via entity extraction + semantic linking) and retrieving evidence with a two-stage process—local entity activation followed by global importance aggregation—yielding stronger and more reliable passage retrieval on multi-hop QA benchmarks. Our implementation is based on \url{https://github.com/DEEP-PolyU/LinearRAG}
    \item \textbf{RAPTOR}~\cite{sarthi2024raptor}: A retrieval-augmented approach that builds a hierarchical tree of recursive embeddings, clusters, and bottom-up summaries, enabling inference-time retrieval across long documents at multiple abstraction levels and delivering strong gains. Our implementation is based on \url{https://github.com/parthsarthi03/raptor}
    \item \textbf{GraphRAG}~\cite{edge2024local}:A graph-based QA framework for private corpora that tackles global, corpus-level questions by (1) building an entity knowledge graph and precomputing community summaries, then (2) answering queries via summary-to-partial-response generation followed by a final aggregation, improving comprehensiveness and diversity over standard RAG at million-token scale. ur implementation is based on \url{https://microsoft.github.io/graphrag/}
\end{itemize}
For all the GraphRAG, we utilize the context of each question as the document and organize the corpus by the official settings. For retrieval, we set the top-k as top-5.
\section{Details of Evaluation Metrics}
\label{app:eval_metric}
\noindent\textbf{Contain Exact Match}~\cite{zhuang2025linearrag}: Check if the correct answer appears in the generated response
\[
\mathrm{Contain\ Exact\ Match}=\frac{1}{N}\sum_{i=1}^{N}\mathbf{1}\big(\mathrm{norm}(a_i)\subseteq \mathrm{norm}(\hat{y}_i)\big)
\]
\noindent\textbf{F-1}: The F1 score evaluates the token-level overlap between the predicted answer $y_i$ and the ground-truth answer $y_i^\star$ using the harmonic mean of precision and recall
\[
\mathrm{F1}=\frac{1}{N}\sum_{i=1}^{N}
\frac{2\cdot\left|\mathrm{tokens}(y_i)\cap \mathrm{tokens}(y_i^\star)\right|}
{\left|\mathrm{tokens}(y_i)\right|+\left|\mathrm{tokens}(y_i^\star)\right|}
\]
\section{F-1 Score for Agentic Search Systems}
In this section, we report the F-1 score of different agentic systems. As Table~\ref{tab:f-1} indicates, both agentic systems can narrow the performance gap between GraphRAG and dense RAG in multi-hop QA settings. For general QA settings, dense-RAG is still a competitive retrieve backend.
\begin{table*}[t]
\centering
\small
\setlength{\tabcolsep}{2pt}
\renewcommand{\arraystretch}{1.10}
\caption{Overall F-1 of  different systems}
\label{tab:f-1}

% Keep: NQ, TriviaQA, HotpotQA, Musique
\begin{tabular}{@{} L{2.2cm} L{3.8cm} *{4}{C{1.5cm}} @{}}
\toprule
\multirow{2}{*}{\textbf{System}} &
\multirow{2}{*}{\textbf{Methods}} &
\multicolumn{2}{c}{\textbf{General QA}} &
\multicolumn{2}{c}{\textbf{Multi-Hop QA}} \\
\cmidrule(r){3-4}\cmidrule(l){5-6}
& & \textbf{NQ$^\dagger$} & \textbf{TriviaQA$^\star$}
  & \textbf{HotpotQA$^\dagger$} & \textbf{Musique$^\star$} \\
\midrule

\multirow{2}{*}{\textbf{Single-shot}} &
Qwen-2.5-7B-Dense & 39.3 & 61.12 & 20.12 & 29.08 \\
& Qwen-2.5-7B-GraphRAG$^\spadesuit$ &
27.92\downdiff{11.38} & 49.25\downdiff{11.87} &
33.72\updiff{13.60} & 41.13\updiff{12.05} \\
\midrule

\multirow{4}{*}{\textbf{Training-free}} &
Search-o1-7B-Dense & 36.53 & 59.73 & 39.08 & 17.46 \\
& Search-o1-7B-GraphRAG$^\spadesuit$ &
36.62\updiff{0.09} & 58.18\downdiff{1.55} &
41.57\updiff{2.49} & 37.01\updiff{19.55} \\
& GraphSearch-7B-Dense & 8.70 & 13.12 & 6.02 & 2.36 \\
& GraphSearch-7B-GraphRAG$^\spadesuit$ &
4.61\downdiff{4.09} & 14.18\updiff{1.06} &
5.48\downdiff{0.54} & 4.21\updiff{1.85} \\
\midrule

\multirow{2}{*}{\textbf{RL-based}} &
Search-R1-7B & 47.26 & 59.21 & 37.13 & 16.02 \\
& Graph-R1-7B$^\spadesuit$ &
44.21\downdiff{3.05} & 62.10\updiff{2.89} &
43.25\updiff{5.12} & 35.12\updiff{19.10} \\
\bottomrule
\end{tabular}
\end{table*}


\section{Efficiency of GraphRAGs}
We report the knowledge construction cost and offline inference time of different graph-based retriever. 
\begin{table*}[!h]
\caption{Comparison on cost of GraphRAGs of NQ. TM indicates construction time per 1M token; CM is the cost per 1M token, RT is the average retrieval time and CT is the average context length.}
\label{tab:cost}
\begin{tabular}{c|cccc}
\hline
\multirow{2}{*}{Method} & \multicolumn{2}{c}{Knowledge Construction} & \multicolumn{2}{c}{Offline Inference} \\ \cline{2-5} 
                        & TM       & \multicolumn{1}{c|}{CM}         & \multicolumn{1}{c|}{RT/s}  & CT/token \\ \hline
HypergraphRAG           & 1.37h    & \multicolumn{1}{c|}{3.93\$}     & \multicolumn{1}{c|}{0.77}  & 1680     \\
HippoRAG2               & 1.19h    & \multicolumn{1}{c|}{2.85\$}     & \multicolumn{1}{c|}{1.00}     & 3229     \\
LinearRAG               & 0.68h    & \multicolumn{1}{c|}{0\$}        & \multicolumn{1}{c|}{1.18}  & 4600     \\
RAPTOR                  & 1.70h    & \multicolumn{1}{c|}{6.38\$}     & \multicolumn{1}{c|}{8.4}   & 814      \\
GraphRAG                & 1.72h    & \multicolumn{1}{c|}{13.19\$}    & \multicolumn{1}{c|}{1.16}  & 22160    \\ \hline
\end{tabular}
\end{table*}
\section{Case Study}
In this section, we displays the difference between different GraphRAGs.
% preamble

% ===== table =====
\begin{table*}[htbp]
\centering
\caption{Case study of different agentic systems.}
\label{tab:case}
\setlength{\tabcolsep}{10pt}       % padding
\renewcommand{\arraystretch}{1.35} % default (will be overridden per-section)

{\scriptsize % <<< 统一整张表字号（可改成 \footnotesize 或 \small）
\begin{tabular}{|>{\centering\arraybackslash}m{1.5cm}|
                    >{\centering\arraybackslash}m{4cm}|
                    >{\centering\arraybackslash}m{4cm}|
                    >{\centering\arraybackslash}m{4cm}|}
\hline

% ---- compact header rows (half height), gray background ----
\noalign{\renewcommand{\arraystretch}{0.675}}
\cellcolor{gray!20}\textbf{Question} &
\multicolumn{3}{>{\columncolor{gray!20}\centering\arraybackslash}m{\dimexpr 12cm + 4\tabcolsep\relax}|}
{\textbf{"Are the Laleli Mosque and Esma Sultan Mansion located in the same neighborhood?"}}\\
\hline

\cellcolor{gray!20}\textbf{Golden Answer} &
\multicolumn{3}{>{\columncolor{gray!20}\centering\arraybackslash}m{\dimexpr 12cm + 4\tabcolsep\relax}|}
{{\textbf{"No"}}}\\
\hline

\cellcolor{gray!20}\textbf{qwen2.5-7B-Ins} &
\cellcolor{gray!20}\textbf{Search-o1-Dense} &
\cellcolor{gray!20}\textbf{Search-R1} &
\cellcolor{gray!20}\textbf{Graph-R1-HippoRAG2} \\
\hline

% ---- body row: fixed-height cells (6cm), vertically centered ----
\noalign{\renewcommand{\arraystretch}{1.35}}
\cellcolor{gray!20}\parbox[t][6cm][c]{1.5cm}{\centering\textbf{Generation}} &
\parbox[t][6cm][c]{4cm}{\centering
To answer the question of whether the Laleli Mosque and Esma Sultan Mansion are located in the same neighborhood, I need to find the locations of both structures.Assistant: \textcolor{orange}{\texttt{<|begin\_search\_query|>}} \\ location of Laleli Mosque
\textcolor{orange}{\texttt{<|end\_search\_query|>}}\\ \textcolor{red}{\texttt{<|begin\_search\_result|>}} The Laleli Mosque is located in Laleli, Fatih, Istanbul, Turkey...\textcolor{red}{\texttt{<|end\_search\_result|>}}\\
Assistant: Based on the information gathered...\{No\}

} &
\parbox[t][6cm][c]{4cm}{\centering
\textcolor{blue}{\texttt{<think>}}To answer the question ``Are the Laleli Mosque and Esma Sultan Mansion located in the same neighborhood?'',
I will first reason through what I know about these two landmarks.
I know that Laleli Mosque is a significant historical site in Istanbul, Turkey.
I need more information about the location of Esma Sultan Mansion to determine if it shares a neighborhood with Laleli Mosque.

\textcolor{blue}{\texttt{</think>}}
\textcolor{orange}{\texttt{<search>}} ``Esma Sultan Mansion location Istanbul Turkey'' \textcolor{orange}{\texttt{</search>}}

\textcolor{red}{\texttt{<information>}} >Doc 1(Title: ``Esma Sultan Mansion''):\ldots\ \textcolor{red}{\texttt{</information>}}
\textcolor{blue}{\texttt{<think>}}\ldots

Based on my search results: Laleli Mosque is in Laleli neighborhood. Esma Sultan Mansion is in Ortak\"oy neighborhood.

Since Laleli Mosque is in Laleli neighborhood and Esma Sultan Mansion is in Ortak\"oy neighborhood,
they are not located in the same neighborhood.
\textcolor{blue}{\texttt{</think>}}

\textcolor{green}{\texttt{<answer>}} No \textcolor{green}{\texttt{</answer>}}
} &
\parbox[t][6cm][c]{4cm}{\centering
\textcolor{blue}{\texttt{<think>}}
To answer this question, I need to determine the neighborhoods where the Laleli Mosque and Esma Sultan Mansion are located.
I will first search for the location of the Laleli Mosque and then for the Esma Sultan Mansion, and finally compare their neighborhoods.
\textcolor{blue}{\texttt{</think>}}

\textcolor{orange}{\texttt{<query>}}
\texttt{\{"query": "Laleli Mosque neighborhood"\}}
\textcolor{orange}{\texttt{</query>}}

\textcolor{red}{\texttt{<knowledge>}}
Wikipedia Title: Laleli Mosque\ldots
\textcolor{red}{\texttt{</knowledge>}}

\textcolor{blue}{\texttt{<think>}}
According to the Wikipedia information retrieved, the Laleli Mosque is located in Laleli, Fatih, Istanbul.
Now, I need to find the neighborhood where Esma Sultan Mansion is located.\\
\textcolor{blue}{\texttt{</think>}}
\ldots

\textcolor{green}{\texttt{<answer>}} No \textcolor{green}{\texttt{</answer>}}
} \\
\hline

\end{tabular}
} % <<< end scriptsize
\end{table*}







\section{Templates}
\subsection{Instructions for Search-o1}
\subsubsection{Instruction for Search-o1}
\begin{figure*}[t]
\centering
\begin{tcolorbox}[mygreenbox, width=\textwidth, title={Instruction for Search-o1},    fonttitle=\Large\bfseries]
You are a reasoning assistant with the ability to perform web searches to help you answer the user’s question
accurately. \\
You have special tools:\\
To perform a search: write \texttt{\textless|begin\_search\_query|\textgreater} your query here \texttt{\textless|end\_search\_query|\textgreater}.\\
Then, the system will search and analyze relevant web pages, then provide you with helpful information in the
format \texttt{\textless|begin\_search\_result|\textgreater}
...search results... \texttt{\textless|end\_search\_result|\textgreater}.\\
You can repeat the search process multiple times if necessary. The maximum number of search attempts is
limited to \texttt{\{MAX\_SEARCH\_LIMIT\}}.\\
Once you have all the information you need, continue your reasoning.\\
Example:\\
Question: ``...''\\
Assistant thinking steps:\\
- I might need to look up details about ...\\
Assistant:\\
\texttt{\textless|begin\_search\_query|\textgreater}...\texttt{\textless|end\_search\_query|\textgreater}
(System returns processed information from relevant web pages)\\
Assistant continues reasoning with the new information...\\
Remember:\\
- Use \texttt{\textless|begin\_search\_query|\textgreater} to request a web search and end with \texttt{\textless|end\_search\_query|\textgreater}.\\
- When done searching, continue your reasoning.\\
Question:\{question\}
\end{tcolorbox}
\caption{Instruction for Search-o1}
\end{figure*}

\subsubsection{Instruction for Reason-in-Documents}
\begin{figure*}[t]
\centering
\begin{tcolorbox}[mygreenbox, width=\textwidth, title={Instruction for Reason-in-Documents}, fonttitle=\Large\bfseries]
Task Instruction:\\
You are tasked with reading and analyzing web pages based on the following inputs: Previous Reasoning Steps, Current Search Query, and Searched Web Pages. Your objective is to extract relevant and helpful information for Current Search Query from the Searched Web Pages and seamlessly integrate this information into the Previous Reasoning Steps to continue reasoning for the original question.\\
Guidelines:\\
1. Analyze the Searched Web Pages:\\
- Carefully review the content of each searched web page.\\
- Identify factual information that is relevant to the Current Search Query and can aid in the reasoning process for the original question.\\
2. Extract Relevant Information:\\
-Select the information from the Searched Web Pages that directly contributes to advancing the Previous Reasoning Steps.\\
-Ensure that the extracted information is accurate and relevant.\\
3. Output Format:\\
- If the web pages provide helpful information for current search query: Present the information beginning with `Final Information' as shown below.\\
Final Information\\
{[Helpful information]}\\
- If the web pages do not provide any helpful information for current search query: Output the following text.\\
Final Information\\
No helpful information found.
Inputs:\\
- Previous Reasoning Steps: \\
\texttt{\{prev\_reasoning\}}\\
- Current Search Query: \\
\texttt{\{search\_query\}}\\
- Searched Web Pages: \\
\texttt{\{document\}}\\
Now you should analyze each web page and find helpful information based on the current search query ``\texttt{\{search\_query\}}'' and previous reasoning steps.
\end{tcolorbox}
\caption{Instruction for Reason-in-Documents}
\end{figure*}

\subsubsection{Instruction for Open-Domain QA Tasks}
\begin{figure*}[h]
\centering
\begin{tcolorbox}[mygreenbox, width=\textwidth, title={Instruction for Open-Domain QA Tasks}, fonttitle=\Large\bfseries]
Please answer the following question.\\
You should provide your final answer in the format \texttt{\textbackslash boxed\{YOUR\_ANSWER\}}.\\
Question:\\
\texttt{\{question\}}
\end{tcolorbox}
\caption{Instruction for Open-Domain QA Tasks}
\end{figure*}

\subsubsection{Additional Notes}

\paragraph{\textbf{Prompting Details}}
For all the instructions above, we input them as user prompts, not system prompts. For non-reasoning models like Qwen2.5-32B-Instruct and Qwen2.5-7B-Instruct, we add a Chain-of-Thought prompt ``You should think step by step to solve it.'' before the question to explicitly make these models reason before giving the final answer.

\paragraph{\noindent\textbf{Implementation Note}}
While the instructions prompt the model to perform ``web searches'', our experiments replaced the actual Bing Web Search API with a retrieval server. The model's generated search queries were intercepted and redirected to our retrieval corpus. The retrieved knowledge chunks were then formatted to mimic web search results before being returned to the model.\\

\subsection{GraphSearch Instruction}
\subsubsection{Query Decomposition}
\begin{figure*}[t]
\centering
\begin{tcolorbox}[mybluebox, width=\textwidth, title={Query Decomposition}]
---Role---\\
You are a helpful assistant specializing in complex query decomposition.\\[2pt]
---Goal---\\
Given a main query, your task is to break it down into several atomic sub-queries, which should
directly correspond to parts of the original query.\\[2pt]
---Instructions---\\
- Decompose the main query into clear and actionable sub-queries that represent smaller, solvable
pieces of the main question.\\
- Ensure that each sub-query addresses one specific entity or concept, with the goal of retrieving
information that will answer the overall main query.\\
- Use sequential numbering (i.e., \texttt{\#1}, \texttt{\#2}, etc.) to represent answers of previous sub-queries. For
example, \texttt{\#1} refers to the answer of Sub-query 1.\\
- Make sure the sub-queries are logically ordered, where the output of one sub-query might feed into
the next.\\
- The final output should be in JSON format, where each sub-query is listed as a key-value pair.\\[2pt]
---Examples---\\
Main Query:\\
How many times did plague occur in the place where the creator of The Worship of Venus died?\\
Sub-queries:\\
\{\{\\
"Sub-query 1": "Who is the creator of The Worship of Venus?",\\
"Sub-query 2": "Where did \texttt{\#1} die?",\\
"Sub-query 3": "How many times did plague occur in \texttt{\#2}?"\\
\}\}\\[4pt]
Main Query:\\
When did the city where Hillcrest High School is located become the capital of the state where the
screenwriter of The Poor Boob was born?\\
Sub-queries:\\
\{\{\\
"Sub-query 1": "Where is Hillcrest High School located?",\\
"Sub-query 2": "Who is the screenwriter of The Poor Boob?",\\
"Sub-query 3": "Where was \texttt{\#2} born?",\\
"Sub-query 4": "When did the city from \texttt{\#1} become the capital of the state from \texttt{\#3}?"\\
\}\}\\[4pt]
Main Query:\\
What crop, which is a big feeder of nitrogen, has a gross income of \$1,363.00 per acre and a net profit
of \$658.00?\\
Sub-queries:\\
\{\{\\
"Sub-query 1": "Which crops are considered big feeders of nitrogen?",\\
"Sub-query 2": "Among \texttt{\#1}, which crop has a gross income of \$1,363.00 per acre?",\\
"Sub-query 3": "Does \texttt{\#2} have a net profit of \$658.00?"\\
\}\}\\[4pt]
---Input---\\
Main Query:\\
\{query\}\\[2pt]
---Output---

\end{tcolorbox}
\caption{GraphSearch Query Decomposition}
\end{figure*}

\subsubsection{Query Decomposition (Knowledge Graph)}
\begin{figure*}[t]
\centering
\begin{tcolorbox}[mybluebox, width=\textwidth, title={Query Decomposition (Knowledge Graph)}]
---Role---\\
You are a helpful assistant specializing in complex query decomposition for knowledge graph retrieval.\\[2pt]
---Goal---\\
Given a main query, your task is to break it down into atomic sub-queries in the form of subject-predicate-object triples. These should correspond directly to parts of the original query and be suitable for querying a knowledge graph.\\[2pt]
---Instructions---\\
- Decompose the main query into a sequence of sub-queries, where each sub-query consists of one or more atomic triples in the format: ("entity1", "relationship", "entity2").\\
- Replace any unknown entity with a placeholder such as Entity\texttt{\#1}, Entity\texttt{\#2}, etc.\\
- Maintain logical ordering, where the result of one sub-query (e.g., Entity\texttt{\#1}) might be required for the next.\\
- Each sub-query may contain more than one triple if needed to express the full meaning.\\
- The final output should be in JSON format, where each key is a sub-query and the value is a list of atomic triples enclosed in parentheses.\\[2pt]
---Examples---\\
Main Query:\\
How many times did plague occur in the place where the creator of The Worship of Venus died?\\
Sub-queries:\\
\{\{\\
"Sub-query 1": [("The Worship of Venus", "is created by", "Entity\texttt{\#1}")],\\
"Sub-query 2": [("Entity\texttt{\#1}", "died at", "Entity\texttt{\#2}")],\\
"Sub-query 3": [\\
("Plague", "occur in", "Entity\texttt{\#2}"),\\
("Plague", "times of occur", "Entity\texttt{\#3}")\\
]\\
\}\}\\[4pt]
Main Query:\\
When did the city where Hillcrest High School is located become the capital of the state where the screenwriter of The Poor Boob was born?\\
Sub-queries:\\
\{\{\\
"Sub-query 1": [("Hillcrest High School", "is located in", "Entity\texttt{\#1}")],\\
"Sub-query 2": [("The Poor Boob", "has screenwriter", "Entity\texttt{\#2}")],\\
"Sub-query 3": [("Entity\texttt{\#2}", "was born in", "Entity\texttt{\#3}")],\\
"Sub-query 4": [\\
("Entity\texttt{\#1}", "is capital of", "Entity\texttt{\#3}"),\\
("Entity\texttt{\#1}", "became capital at", "Entity\texttt{\#4}")\\
]\\
\}\}\\[4pt]
Main Query:\\
What crop, which is a big feeder of nitrogen, has a gross income of \$1,363.00 per acre and a net profit of \$658.00?\\
Sub-queries:\\
\{\{\\
"Sub-query 1": [("Entity\texttt{\#1}", "is a", "crop that is a heavy nitrogen feeder")],\\
"Sub-query 2": [("Entity\texttt{\#1}", "has gross income per acre", "\$1,363.00")],\\
"Sub-query 3": [("Entity\texttt{\#1}", "has net profit", "\$658.00")]\\
\}\}\\[4pt]
---Input---\\
Main Query:\\
\{query\}\\[2pt]
---Output---
\end{tcolorbox}
\caption{GraphSearch Query Decomposition (KG)}
\end{figure*}

\subsubsection{Evidence Verification}
\begin{figure*}[t]
\centering
\begin{tcolorbox}[mybluebox, width=\textwidth, title={Evidence Verification}]
---Role---\\
You are a critical evaluator specializing in verifying the logical soundness and evidential sufficiency of model-generated responses.\\[2pt]
---Goal---\\
Given a user query, retrieved context data, and the model-generated response, your task is to evaluate whether the response forms a rigorous
logical loop supported by the provided evidence.\\[2pt]
---Instructions---\\
- Carefully examine whether the response is \textbf{strictly grounded} in the retrieved context data.\\
- Assess whether the reasoning process forms a \textbf{complete logical chain}, without missing steps or unsupported leaps.\\
- Identify if there are \textbf{evidence gaps, low-confidence claims, or speculative statements}.\\
- If the response demonstrates a well-supported, confident, and logically closed argument, conclude your analysis with \textbf{``Yes''}.\\
- If the response shows hesitation, incomplete reasoning, or lacks solid evidence support, conclude your analysis with \textbf{``No''}.\\[2pt]
---Input---\\
User-Query:\\
\{query\}\\
Retrieved Context Data:\\
\{context\_data\}\\
Model Response:\\
\{model\_response\}\\[2pt]
---Output---
\end{tcolorbox}
\caption{GraphSearch Evidence Verification}
\end{figure*}

\subsubsection{Deep Answer Generation}
\begin{figure*}[t]
\centering
\begin{tcolorbox}[mybluebox, width=\textwidth, title={Deep Answer Generation}]
---Role---\\
You are a helpful assistant specializing in complex question answering.\\[2pt]
---Goal---\\
Given a complex query and retrieved context data, your task is to construct a logically sound, step-by-step answer. 
Your explanation should follow a rigorous reasoning path, incorporate relevant evidence, and establish clear relationships between the entities.\\[2pt]
---Instructions---\\
- Break down the reasoning process into clear, coherent steps.\\
- Use context data explicitly to support each reasoning step.\\
- Make sure relationships between entities are logically explained.\\[2pt]
---Input---\\
Query:\\
\{query\}\\
Context Data:\\
\{context\_data\}\\[2pt]
---Output---
\end{tcolorbox}
\caption{GraphSearch Deep Answer Generation}
\end{figure*}

\subsubsection{Query Expansion}
\begin{figure*}[t]
\centering
\begin{tcolorbox}[mybluebox, width=\textwidth, title={Query Expansion}]
---Role---\\
You are a helpful assistant specializing in query expansion for evidence completion.\\[2pt]
---Goal---\\
Given a main query, retrieved context data, the model-generated response, and the evidence verification analysis, your task is to perform \textbf{query expansion}.\\
If the evidence verification analysis shows that the current evidence is insufficient to support the logical chain of the response, generate one or more additional sub-queries.\\
These sub-queries should aim to cover missing retrieval scenarios, fill in the evidence gaps, and guide towards a more complete and confident logical reasoning chain.\\[2pt]
---Instructions---\\
- Use the retrieved context data, especially any existing sub-queries in the retrieval history, as references when generating new sub-queries.\\
- Focus on producing \textbf{complementary sub-queries} that address aspects not yet fully supported by evidence.\\
- Avoid duplicating existing sub-queries; instead, expand into related but uncovered areas.\\
- Keep sub-queries clear, specific, and directly actionable for retrieval.\\
- Output should be in the form of a \textbf{Python-style List of strings}, where each string is a new sub-query.\\[2pt]
---Input---\\
Main Query:\\
\{query\}\\
Retrieved Context Data:\\
\{context\_data\}\\
Model Response:\\
\{model\_response\}\\
Evidence Verification Analysis:\\
\{evidence\_verification\}\\[2pt]
---Output---
\end{tcolorbox}
\caption{GraphSearch Query Expansion}
\end{figure*}

\subsection{Search-R1 Instruction}
\begin{figure*}[t]
\centering
\begin{tcolorbox}[mygraybox, width=\textwidth, title={Search-R1 Instruction}]
Answer the given question.

You must conduct reasoning inside <think> and </think> first every time you get new information.

After reasoning, if you find you lack some knowledge, you can call a search engine by <search> query </search> and it will return the top searched results between <information> and </information>.

You can search as many times as your want.

If you find no further external knowledge needed, you can directly provide the answer inside <answer> and </answer>, without detailed illustrations. For example, <answer> Beijing </answer>. 

Question: \{question\}
\end{tcolorbox}
\caption{Search-R1 Instruction}
\end{figure*}

\subsection{Graph-R1 Instruction}
\begin{figure*}[t]
\centering
\begin{tcolorbox}[mypurplebox, width=\textwidth, title={Graph-R1 Instruction},    fonttitle=\Large\bfseries]
Answer the given question. \\
You can query from knowledge base provided to you to answer the question. \\
You can query knowledge as many times as you want.
You must first conduct reasoning inside <think>...</think>. \\
If you need to query knowledge, you can set a query statement between <query>...</query> to query from knowledge base after <think>...</think>. \\
When you have the final answer, you can output the answer inside <answer>...</answer>. \\
Output format for query: \\
<think>\\
... \\
</think>\\
<query>\\
... \\
</query> \\
Output format for answer:\\
<think>\\
... \\
</think>\\
<answer>\\
... \\
</answer> \\
Question: \{question\}
\end{tcolorbox}
\caption{Graph-R1 Instruction}
\end{figure*}